\documentclass{article}
\usepackage{graphicx}
\usepackage{float} 

\title{Tractive System Pre-charger Module (First Draft)}
\author{William Bowley (S4226716) | Recruit Power Train}
\date{\today}

\begin{document}
\maketitle

\section{Overview}
The pre-charge module exists to mitigate inrush current when the accumulator is connected to the 
tractive system. This is because of the existence of a large bulk capacitance placed in parallel 
with the accumulator to manage line transients from large dynamic current sinks like the motors.

\vspace{0.5cm}
The high-level system architecture was predefined by the system leads to be based on a passive RC 
(resistor \& capacitor) circuit because start-up isn't a competition criterion and hence increasing 
$\tau$ to expand the duration of the current pulse is optimal and it increases margins. Most importantly, 
the passive pre-charger is extremely reliable due to minimal failure points compared to alternatives.


\vspace{0.5cm}
The pre-charge sequence is as follows: the negative HV contactor connects the accumulator negative terminal to the
system reference. The pre-charge contactor then is allowed to close, which enables current flow through the 
pre-charge resistor into the DC link capacitance. Once the DC link voltage approaches the accumulator voltage
approximately 3$\tau$, the main positive contactor is closed. The pre-charge contactor is subsequently opened
to remove the series resistance from the power path. 

\section{Constraints}

Solid-state relays or FETs are not allowed due to RE.5.6.5 `The pre-charge relay must be a mechanical type relay'
hence a mechanical relay will be used as the pre-charge contactor. All components that dissipate heat must be 
monitored for thermal overload. The PDOC must open the shutdown circuit before exceeding the maximum operating temperature.
However, due to this design being validated to run continuously even in a faulted state, the PDOC can be omitted.

\vspace{0.5cm}
Importantly, the mechanical relay must be normally open such that if shutdown (SDC) is low, the accumulator must be disconnected
from the tractive system. 

\section{Detailed Design}
The PWR263S-35 power resistor was chosen due to its 35 W rating and its package being suitable for SMD (TO-263). It was decided that
using a few resistors in series would be better as they distribute the heat. The KT24-1A-40L-SMD mechanical switch was
chosen because of its 1 A continuous current rating and low coil resistance of 150 to 200 m$\Omega$.

\vspace{0.5cm}
The activation time was arbitrarily set to 4.5 s which at a 3$\tau$ requirement resulted in 6x1 k$\Omega$ resistors from `system\_stress.py'.
The maximal inrush current hence was $\sim$100 mA with 43.2 J being delivered to the DC link capacitor.

\vspace{0.5cm}
A worst case continuous fault leads to each resistor absorbing 10 W continuous which is $\sim$1/3.5 of their maximum rating.
The mechanical relay is 
connected to \texttt{pre\_aux\_out} \& \texttt{pre\_aux\_in} from the AMS main board to enable or disable it. With \texttt{prech\_coil\_pos} \&
\texttt{prech\_coil\_neg} supplying the current to activate the relay. This design reflects the R26 AMS architecture. 

\vspace{0.5cm}
A 4 pin XH 2.54 mm JST connector was selected due to it being the standard for a low voltage application like the AMS connection to pre-charger. A Molex
connector was used for the input and output of the HV line specifically a Molex Mini-Fit Jr which is rated for 600 V while having a reasonable form factor.
However it is a 4 pin connector hence only 2 pins are used which are orthogonal to each other. The trace width is set to 0.80 mm and the trace spacing is
1.25 mm for both low voltage and high voltage components.


\section{Conclusion}
The design proposed here should function correctly under standard operating conditions. Images and data-sheets can be found within `1-deliverable-system'.
The Bourns PWR263S-35-1001F resistor should have a significant safety margin for both power dissipation (10 W vs 35 W) and operating voltage (600 V vs 700 V).

By designing a series-string passive pre-charger, the thermal load is distributed across the PCB. Furthermore, the integration of a mechanical reed relay 
ensures full compliance with Rule RE.5.6.5, providing the required isolated fail-safe state when the shutdown circuit is de-energized. This approach
balances the robust and reliable requirements of a tractive system and also the strict safety constraints.


\section{Bill of Materials}
\begin{table}[H]
\centering
\caption{Estimated Component and Manufacturing Costs (AUD)}
\begin{tabular}{|l|l|c|c|r|}
\hline
    \textbf{Component} & \textbf{Part Number} & \textbf{Qty} & \textbf{Unit Price} & \textbf{Total} \\ \hline
    Power Resistor (1k$\Omega$) & Bourns PWR263S-35-1001F & 6 & \$5.15 & \$30.90 \\ \hline
    Reed Relay & Standex KT24-1A-40L-SMD & 1 & \$8.40 & \$8.40 \\ \hline
    HV Header (4-pin) & Molex Mini-Fit Jr (39-29-1048) & 2 & \$1.70 & \$3.40 \\ \hline
    LV Header (4-pin) & JST XH (B4B-XH-A) & 1 & \$0.21 & \$0.21 \\ \hline
    PCB Manufacturing & Custom 2-Layer Board & 1 & \$10.00 & \$10.00 \\ \hline
    \multicolumn{4}{|l|}{\textbf{Subtotal}} & \textbf{\$52.91} \\ \hline
\end{tabular}
\end{table}

\end{document}